\documentclass[11pt,a4paper]{article}

\usepackage[utf8]{inputenc}
\usepackage{amsmath,amssymb,amsthm,mathtools}
\usepackage{booktabs,array}
\usepackage[margin=2.5cm]{geometry}
\usepackage{xcolor}
\usepackage[colorlinks=true,linkcolor=blue!50!black,citecolor=blue!50!black,urlcolor=blue!50!black]{hyperref}

% ---------- theorem environments ----------
\newtheorem{theorem}{Theorem}[section]
\newtheorem{proposition}[theorem]{Proposition}
\theoremstyle{definition}
\newtheorem{example}[theorem]{Example}
\theoremstyle{remark}
\newtheorem{remark}[theorem]{Remark}

% ---------- macros ----------
\newcommand{\dd}{\mathrm{d}}
\newcommand{\R}{\mathbb{R}}
\newcommand{\Tor}{\mathbb{T}}
\newcommand{\PB}[2]{\{#1,#2\}}
\newcommand{\weq}{\approx}
\newcommand{\FL}{\mathbb{F}\!L}
\newcommand{\Sw}{\mathcal{S}}
\DeclareMathOperator{\rank}{rank}

\title{Coordinates and Momenta Are Not Equal Partners:\\
Action--Angle Variables, Dirac's Constraints and Singular Lagrangians}
\author{}
\date{October 2026}

\begin{document}
\maketitle

\begin{abstract}
\noindent
The textbook statement that in phase space ``coordinates and momenta are on an equal footing'' is exact only for the symplectic form, that is, for the \emph{kinematics} of Hamiltonian mechanics. Once the data that turn a Hamiltonian system into a mechanical system are added --- the Legendre fibration $T^*Q\to Q$, a Hamiltonian with its integrals of motion, and the algebra of constraints --- coordinates and momenta play different roles. We discuss three manifestations of this asymmetry and show that they are aspects of one structure.
(i)~In action--angle variables the momenta (actions) are constants of motion while the conjugate angles evolve; because the angles are periodic, no global canonical transformation can exchange the two.
(ii)~For a singular Lagrangian the Legendre map is a bundle map over the configuration space, so primary constraints are always ``vertical'': they restrict momenta and never coordinates alone.
(iii)~In Dirac's classification, in an adapted Darboux chart, first-class constraints are momenta whose conjugate coordinates are pure gauge, whereas second-class constraints come in conjugate $(q,p)$ pairs and are, together with the first-order action, where the symmetric picture survives.
We conclude that the invariant distinction is not ``$q$ versus $p$'' but ``generator of a symmetry versus parameter along its orbit''.
\end{abstract}

%=====================================================================
\section{Introduction}
%=====================================================================

It is a commonplace of analytical mechanics that the Hamiltonian formalism puts coordinates and momenta ``on an equal footing''. Three facts are usually quoted in support. First, canonical transformations are allowed to mix $q$ and $p$ freely; the exchange $(q,p)\mapsto(p,-q)$ is canonical. Second, Hamilton's equations $\dot q=\partial H/\partial p$, $\dot p=-\partial H/\partial q$ differ only by a sign. Third, in quantum mechanics the position and momentum representations are related by a Fourier transform and are equally good.

Elementary experience in mechanics nevertheless suggests that the two kinds of variables are not interchangeable.
For a free particle $p$ is constant and $q$ is not.
For an integrable system with compact invariant sets there are canonical variables $(\theta,I)$ in which the \emph{momenta} $I$ (the action variables) are constants of motion while the \emph{coordinates} $\theta$ (the angles) change linearly in time; the opposite situation --- constant coordinates and evolving momenta --- is never presented as the normal form of regular motion.
In gauge theories some momenta vanish identically while their conjugate coordinates are arbitrary functions of time.
In the Lagrangian formalism the passage to momenta is the Legendre map $\dot q\mapsto p=\partial L/\partial\dot q$, which fails to be invertible exactly when the Lagrangian is singular.

The purpose of this note is to show that these observations are not accidents but manifestations of a single fact: \emph{the folklore statement is true of the symplectic form, but a mechanical system carries structure beyond the symplectic form, and with respect to that structure $q$ and $p$ are inequivalent.}
Our results, in the order in which they appear, are:
\begin{itemize}
\item[\textbf{R1}] The canonical transformations compatible with Lagrangian mechanics on a configuration space $Q$ form the small group $\mathrm{Diff}(Q)\ltimes Z^1(Q)$, which does not mix $q$ and $p$ (Proposition~\ref{prop:fibre}).
\item[\textbf{R2}] In action--angle variables the exchange of an action and its angle is locally possible but globally obstructed (Theorem~\ref{thm:noconj}). The action--angle Lagrangian is regular iff the frequency map is a local diffeomorphism (Kolmogorov non-degeneracy).
\item[\textbf{R3}] Primary constraints of a singular Lagrangian are always vertical: no primary constraint involves coordinates alone (Proposition~\ref{prop:vertical}).
\item[\textbf{R4}] In Dirac's classification, first-class constraints are, in an adapted chart, momenta with pure-gauge conjugates, while second-class constraints come in conjugate pairs (Proposition~\ref{prop:normal}).
\item[\textbf{R5}] The first-order action $\int(p\dot q-H)\,\dd t$ \emph{is} symmetric in $q$ and $p$ up to a boundary term; this is the sense in which the folklore is exactly right (Section~\ref{sec:first-order}).
\end{itemize}

\paragraph{Conventions and assumptions.}
$Q$ is an $n$-dimensional smooth manifold, $\pi:T^*Q\to Q$ its cotangent bundle with local coordinates $(q^i,p_i)$, canonical one-form $\vartheta=p_i\,\dd q^i$ and symplectic form $\omega=\dd\vartheta=\dd p_i\wedge\dd q^i$. The Poisson bracket is $\PB fg=\partial_{q^i}f\,\partial_{p_i}g-\partial_{p_i}f\,\partial_{q^i}g$, so $\PB{q^i}{p_j}=\delta^i_j$ and $\dot f=\PB fH$. The flow of a function $f$ is $\dd g/\dd s=\PB gf$. Weak equality (on the constraint surface) is written $\weq$. We work with finitely many degrees of freedom and assume the usual constant-rank and irreducibility conditions of the Dirac--Bergmann theory \cite{gnh,ht}; field-theoretic examples are treated per spatial point.

%=====================================================================
\section{What is true: the symplectic statement}
%=====================================================================

\subsection{The exchange map and its consequences}

The folklore statement is
\begin{quote}
\textbf{(F)} \emph{The variables $q$ and $p$ are equivalent: any canonical pair can be taken as ``coordinate'' and ``momentum'' in either order.}
\end{quote}
Its basis is the map
\begin{equation}\label{eq:swap}
\Sw:(q,p)\longmapsto(Q,P)=(p,-q),\qquad \Sw^*\omega=\dd P_i\wedge\dd Q^i=-\dd q^i\wedge\dd p_i=\omega ,
\end{equation}
which is a symplectomorphism of $\R^{2n}$. More generally, the linear canonical transformations form $\mathrm{Sp}(2n,\R)$, which contains maps that mix $q$'s and $p$'s arbitrarily, and by Darboux's theorem every function with non-vanishing differential is a coordinate of some local canonical chart \cite{arnold,am}. In quantum mechanics on $\R^n$ the map $\Sw$ is implemented by the Fourier transform. Hence for a particle on $\R^n$, without constraints, the \emph{kinematic} structure is exactly symmetric. Statement (F) is correct in this setting, and we do not dispute it there.

What (F) does \emph{not} address is the extra structure that a mechanical system carries. The most basic piece is the \emph{polarisation} $\pi:T^*Q\to Q$ produced by the Legendre map of a Lagrangian $L(q,\dot q)$. Its fibres are the sets of momenta at a given configuration; $q$ labels the base, $p$ the fibre. The following statement makes precise how much of $\mathrm{Symp}(T^*Q)$ respects this structure.

\begin{proposition}[fibre-preserving canonical transformations]\label{prop:fibre}
Let $\Phi:T^*Q\to T^*Q$ be a symplectomorphism mapping every fibre $T^*_qQ$ onto a fibre. Then $\Phi=\hat\varphi\circ\tau_\alpha$, where $\varphi:Q\to Q$ is a diffeomorphism, $\hat\varphi$ is its cotangent lift, and $\tau_\alpha(q,p)=(q,p+\alpha(q))$ with $\alpha$ a \emph{closed} one-form on $Q$. In coordinates, $Q^a=\varphi^a(q)$ and $P_a=\frac{\partial q^i}{\partial Q^a}\,(p_i+\alpha_i(q))$.
\end{proposition}

\begin{proof}
Let $\varphi$ be the diffeomorphism of $Q$ covered by $\Phi$. Cotangent lifts preserve $\vartheta$, hence $\omega$, so $\Psi=\hat\varphi^{-1}\circ\Phi$ is a symplectomorphism that covers the identity: $\Psi(q,p)=(q,\psi(q,p))$. The condition $\Psi^*\omega=\omega$ reads
$\dd\psi_i\wedge\dd q^i=\dd p_i\wedge\dd q^i$. Comparing the coefficients of $\dd p_j\wedge\dd q^i$ gives $\partial\psi_i/\partial p_j=\delta_i^j$, hence $\psi_i=p_i+\alpha_i(q)$. The remaining terms give $(\partial_j\alpha_i-\partial_i\alpha_j)\,\dd q^j\wedge\dd q^i=0$, i.e.\ $\dd\alpha=0$.
\end{proof}

Thus the transformations under which Lagrangian mechanics on $Q$ is covariant --- point transformations $q\mapsto Q(q)$ and gauge shifts $p\mapsto p+\nabla f$ --- form the group $\mathrm{Diff}(Q)\ltimes Z^1(Q)$, a tiny subgroup of the symplectomorphisms. Under it $q$ transforms by an arbitrary diffeomorphism and $p$ as a covector (affinely, up to a closed shift); \emph{no element mixes them}. The exchange $\Sw$ lies outside this group, and it does not preserve the property of being Lagrangian-realisable: a Hamiltonian is \emph{regular} if $\det(\partial^2H/\partial p_i\partial p_j)\neq0$, and this property is invariant under $\mathrm{Diff}(Q)\ltimes Z^1(Q)$ but not under $\Sw$.

\begin{remark}
Time reversal $(q,p)\mapsto(q,-p)$ is anti-symplectic and is a symmetry of every mechanical Hamiltonian that is even in $p$; $q$ is even and $p$ odd under it. Conjugating by $\Sw$ turns it into $(Q,P)\mapsto(-Q,P)$, which is no longer ``reversal of velocities''. The equality (F) does not survive the requirement that the time-reversal structure be kept.
\end{remark}

\subsection{Two examples}

\begin{example}[free particle]\label{ex:free}
For $H=p^2/2m$ the momentum is conserved and $q$ grows linearly. Apply $\Sw$: $\tilde H=Q^2/2m$, with $\dot Q=0$ and $\dot P=-Q/m$. This is a perfectly good Hamiltonian system --- the free particle in disguise --- in which a \emph{coordinate} is conserved. But $\partial^2\tilde H/\partial P^2=0$: the velocity $\dot Q$ vanishes identically while $P$ is arbitrary, so $\tilde H$ is not the Legendre transform of any Lagrangian $L(Q,\dot Q)$. In symplectic geometry both descriptions are equally legitimate; only the first is a \emph{mechanical} one.
\end{example}

\begin{example}[harmonic oscillator]\label{ex:osc}
Take $H=\tfrac12(p^2+q^2)$ and the action--angle chart
\begin{equation}\label{eq:osc-aa}
q=\sqrt{2I}\,\sin\theta,\qquad p=\sqrt{2I}\,\cos\theta,\qquad
\dd p\wedge\dd q=\dd I\wedge\dd\theta,\qquad H=I,\qquad I=\frac1{2\pi}\oint p\,\dd q .
\end{equation}
Then $\PB\theta I=1$, $\dot I=0$, $\dot\theta=1$. The exchange $\Sw$ maps the oscillator to itself, and in the chart~\eqref{eq:osc-aa} it reads
\[
\Sw:(\theta,I)\longmapsto(\theta+\tfrac\pi2,\,I),
\]
because $Q=p=\sqrt{2I}\sin(\theta+\tfrac\pi2)$ and $P=-q=\sqrt{2I}\cos(\theta+\tfrac\pi2)$. So $\Sw$ is the time-$\pi/2$ flow. The perfect $q$--$p$ symmetry of the oscillator is nothing but the statement that the \emph{action} is invariant while the \emph{angle} can be rotated; $\Sw$ exchanges $q$ and $p$ but cannot exchange $I$ and $\theta$. (The chart is singular at the equilibrium $I=0$, like polar coordinates.)
\end{example}

%=====================================================================
\section{Action--angle variables}
%=====================================================================

\subsection{The normal form of regular motion}

\begin{theorem}[Liouville--Arnold \cite{arnold}]\label{thm:LA}
Let $F_1=H,\dots,F_n$ be functions on a $2n$-dimensional symplectic manifold with $\PB{F_j}{F_k}=0$ and $\dd F_1\wedge\dots\wedge\dd F_n\neq0$ on a compact connected level set $M_f=\{F=f\}$. Then $M_f\cong\Tor^n$, and a neighbourhood of $M_f$ carries canonical coordinates $(\theta,I)\in\Tor^n\times U$, $U\subset\R^n$ open, with $\omega=\dd I_k\wedge\dd\theta_k$, such that $I=I(F)$; in particular $H=H(I)$. The actions are $I_k=\frac1{2\pi}\oint_{\gamma_k}p\,\dd q$ for a basis $\gamma_k$ of $H_1(M_f)$.
\end{theorem}

Hamilton's equations become
\begin{equation}\label{eq:aa-eom}
\dot I_k=0,\qquad \dot\theta_k=\nu_k(I):=\frac{\partial H}{\partial I_k},\qquad \theta(t)=\theta_0+\nu(I)\,t .
\end{equation}
Thus $n$ of the $2n$ canonical variables are constants of motion and the other $n$ evolve. This is the general form of regular bounded motion. Locally, the roles can be exchanged: on the universal cover $\R^n\times U$ the map $(\theta,I)\mapsto(Q,P)=(I,-\theta)$ is canonical and gives $H=H(Q)$, $\dot Q=0$, $\dot P=-\nu(Q)$ --- the same pattern as in Example~\ref{ex:free}. What forbids this exchange globally is the topology.

\subsection{The global obstruction}

\begin{theorem}[no global conjugate to a periodic action]\label{thm:noconj}
On $\Tor^n\times U$ with $\PB{\theta_k}{I_l}=\delta_{kl}$ there is no smooth real function $\Theta$ with $\PB\Theta{I_k}=1$ for some $k$. Consequently no canonical transformation defined on all of $\Tor^n\times U$ takes the action $I_k$ to a coordinate.
\end{theorem}

\begin{proof}
The flow of $I_k$ is $\theta_k\mapsto\theta_k+s$, with all other variables fixed, and it is periodic with period $2\pi$. Along it, $\dd\Theta/\dd s=\PB\Theta{I_k}=1$, so $\Theta$ would have to increase by $2\pi$ over a closed orbit --- a contradiction. If $(Q,P)$ were canonical with $Q_k=I_k$, then $\PB{Q_k}{P_k}=1$ gives $\PB{-P_k}{I_k}=1$, and $\Theta=-P_k$ would be such a function.
\end{proof}

\begin{remark}
Topologically, the phase space $T^*\Tor^n=\Tor^n\times\R^n$ carries two Lagrangian fibrations: the vertical one, $\{\theta=\mathrm{const}\}$, with fibres $\R^n$ (the momenta), and the Liouville one, $\{I=\mathrm{const}\}$, with compact fibres $\Tor^n$. These fibres are not homeomorphic, so no symplectomorphism exchanges the two fibrations. The Hamiltonian selects the second fibration: $H$ is constant on its leaves and the motion is tangent to them.
\end{remark}

\subsection{Actions as Noether charges}

The torus acts on $\Tor^n\times U$ by $\theta\mapsto\theta+s$; the action is Hamiltonian with momentum map $I:\Tor^n\times U\to\R^n$. The invariance of $H$ under it is equivalent to $H=H(I)$ and to the conservation of $I$. The actions are therefore Noether charges of the angle translations, formally identical to the conservation of $p_i$ for a cyclic coordinate $q^i$:
\[
\partial H/\partial q^i=0\ \Longrightarrow\ \dot p_i=0 .
\]
The reverse implication --- ``if $H$ does not depend on $p_i$ then $q^i$ is conserved'' --- is formally true ($\dot q^i=\partial H/\partial p_i=0$) but describes a Hamiltonian that is singular in the sense of Example~\ref{ex:free}. More generally, for a vector field $\xi$ on $Q$ that generates a symmetry, the cotangent lift has the conserved charge $J_\xi=p_i\,\xi^i(q)$, which is \emph{linear in the momenta}. A momentum is a \emph{generator}; a coordinate (an angle) is the \emph{parameter} along the orbit of that generator.

\subsection{Link with singular Lagrangians: Kolmogorov non-degeneracy}\label{sec:kolmogorov}

When does the action--angle Hamiltonian $H(I)$ possess a Lagrangian $L(\theta,\dot\theta)$? The Legendre map is $I\mapsto\dot\theta=\nu(I)=\partial H/\partial I$, a local diffeomorphism iff
\begin{equation}\label{eq:kolmogorov}
\det\frac{\partial\nu_k}{\partial I_l}=\det\frac{\partial^2H}{\partial I_k\,\partial I_l}\neq0 ,
\end{equation}
which is Kolmogorov's non-degeneracy condition \cite{arnold}. If~\eqref{eq:kolmogorov} holds then
\[
L=\mathcal L(\dot\theta)=\Bigl[I\cdot\dot\theta-H(I)\Bigr]_{I=\nu^{-1}(\dot\theta)},\qquad I=\frac{\partial\mathcal L}{\partial\dot\theta}.
\]
Hence \emph{every non-degenerate integrable system is, locally, a system on a torus all of whose coordinates are cyclic, and the actions are the Noether momenta of those cyclic coordinates.} If~\eqref{eq:kolmogorov} fails, the action--angle Lagrangian is singular. This happens for the harmonic oscillator, $H=\nu I$, and for the Kepler problem, whose Hamiltonian depends on one linear combination of the actions. In that case there is no $\mathcal L(\dot\theta)$ at all: $\dot\theta=\nu$ identically, and the only available Lagrangian is the first-order $I(\dot\theta-\nu)$, in which $I$ is a Lagrange multiplier. We return to this in Section~\ref{sec:first-order}.

\subsection{Quantum remark}

At the quantum level the same obstruction is the phase problem. The actions are quantised ($I_k=\hbar(n_k+\mu_k/4)$ in the semiclassical EBK scheme, with Maslov indices $\mu_k$), so $\hat I$ has a discrete spectrum, while only the exponentials $e^{\pm i\theta}$ make sense and act as shifts on the lattice of $n$. For the oscillator, with $n\ge0$, the exponential of the phase is not even unitary \cite{dirac27,sg64}. The Fourier duality is therefore between $\ell^2(\mathbb Z)$ and $L^2(S^1)$, not between two copies of $L^2(\R)$.

%=====================================================================
\section{Singular Lagrangians: the Legendre map}
%=====================================================================

Let $L:TQ\to\R$ and let
\[
\FL:TQ\to T^*Q,\qquad \FL(q,\dot q)=\Bigl(q,\ \frac{\partial L}{\partial\dot q}\Bigr)
\]
be the fibre derivative (Legendre map), with Hessian $W_{ij}=\partial^2L/\partial\dot q^i\partial\dot q^j$. The Lagrangian is \emph{regular} if $\det W\neq0$, so that $\FL$ is a local diffeomorphism; then $H(q,p)=p\,\dot q(q,p)-L$ and $H_{pp}=W^{-1}$, so $L$ regular iff $H$ regular. It is \emph{singular} if $\rank W=r<n$. We assume the \emph{almost regular} situation: $\Sigma_1:=\FL(TQ)$ is a submanifold and $\FL:TQ\to\Sigma_1$ is a submersion with connected fibres \cite{gnh}.

\begin{proposition}[primary constraints are vertical]\label{prop:vertical}
Let $L$ be almost regular with $\rank W=r$. Then:
\begin{enumerate}
\item[(i)] $\dim\Sigma_1=n+r$; locally $\Sigma_1$ is cut out by $n-r$ independent primary constraints $\phi_m(q,p)\weq0$.
\item[(ii)] The projection $\pi|_{\Sigma_1}:\Sigma_1\to Q$ is a surjective submersion with fibres of dimension $r$.
\item[(iii)] In particular, no primary constraint is a function of the coordinates alone: if $f(q)$ vanishes on $\Sigma_1$, then $f\equiv0$.
\end{enumerate}
\end{proposition}

\begin{proof}
In the coordinates $(q,\dot q)\mapsto(q,p)$ the differential of $\FL$ is block triangular,
$\begin{psmallmatrix}1&0\\ \ast&W\end{psmallmatrix}$, and has rank $n+r$; this gives (i). Moreover $\FL$ covers the identity: $\pi\circ\FL=\pi_Q$, where $\pi_Q:TQ\to Q$ is a surjective submersion. Hence $\pi|_{\Sigma_1}$ is a surjective submersion, which gives (ii), and its fibres have dimension $\dim\Sigma_1-n=r$. For (iii): if $f(q)$ vanished on $\Sigma_1$, then, since $\pi(\Sigma_1)=Q$, $f$ would vanish on all of $Q$.
\end{proof}

Proposition~\ref{prop:vertical} is the asymmetry in its simplest form. All $n$ coordinates are independent variables; only $r$ of the $n$ momenta are. The constraint surface cuts the fibres of $T^*Q$ but never the base. The exchange $\Sw$ would turn a relation $\phi(q,p)=p_m-f(q,p_a)=0$ into a relation that no longer restricts a fibre and no longer descends from any Lagrangian.

The velocities behave in the complementary way. The equations $p=\partial L/\partial\dot q$ determine only $r$ combinations of the velocities in terms of $(q,p)$; those along the null directions of $W$ remain undetermined, and the Euler--Lagrange equations $W_{ij}\ddot q^j=\dots$ cannot be solved for all accelerations. In the Hamiltonian description the undetermined velocities become the multipliers $u^m$ in
\begin{equation}\label{eq:HT}
H_T=H_c+u^m\phi_m,\qquad \dot q^i=\frac{\partial H_c}{\partial p_i}+u^m\frac{\partial\phi_m}{\partial p_i}.
\end{equation}
\emph{The momentum is constrained; its conjugate velocity is free.}

%=====================================================================
\section{Dirac's constraints}
%=====================================================================

\subsection{The algorithm and the two classifications}

Dirac's procedure \cite{dirac50,dirac64} proceeds in steps. The primary constraints $\phi_m\weq0$ come from the Legendre map. Consistency of $\phi_m\weq0$ with the evolution generated by $H_T$ in~\eqref{eq:HT},
\[
\{\phi_m,H_c\}+u^{m'}\{\phi_m,\phi_{m'}\}\weq0 ,
\]
either fixes some multipliers or produces new, \emph{secondary} constraints; the procedure is iterated until it closes. The final set $\{\chi_A\}$ is classified as:
\begin{itemize}
\item \emph{first class}: $\PB{F}{\chi_A}\weq0$ for all $A$; these generate gauge transformations and leave their multipliers arbitrary;
\item \emph{second class}: the matrix $C_{\alpha\beta}=\PB{\chi_\alpha}{\chi_\beta}$ is invertible on the subset of second-class constraints (which is therefore even in number); the multipliers are fixed by consistency and one passes to the Dirac bracket $\PB FG_D=\PB FG-\PB F{\chi_\alpha}(C^{-1})^{\alpha\beta}\PB{\chi_\beta}G$.
\end{itemize}
With $N_1$ first-class and $N_2$ second-class constraints the reduced phase space has dimension $2n-2N_1-N_2$.

Two different dichotomies thus appear, \emph{primary vs.\ secondary} and \emph{first vs.\ second class}; both display the asymmetry.

\subsection{A local normal form}

\begin{proposition}[adapted chart]\label{prop:normal}
Let $\Sigma\subset T^*Q$ be the final constraint surface, with $k$ first-class and $2m$ second-class constraints, and assume the constant-rank conditions of Section~1. Then near every point of $\Sigma$ there are canonical coordinates
\[
(q^a,p_a;\ Q^\alpha,P_\alpha;\ y^A,\eta_A),\qquad a\le k,\ \alpha\le m,\ A\le n-k-m,
\]
in which $\Sigma=\{p_a=0,\ Q^\alpha=0,\ P_\alpha=0\}$.
\end{proposition}

\begin{proof}[Sketch]
The second-class constraints can be brought, by a symplectic Gram--Schmidt procedure, to canonical pairs $(Q^\alpha,P_\alpha)$. After adding suitable combinations of the second-class constraints, $\phi_a\to\phi_a-\PB{\phi_a}{\chi_\alpha}(C^{-1})^{\alpha\beta}\chi_\beta$, the first-class constraints commute with $Q^\alpha,P_\alpha$ and define a coisotropic submanifold of the symplectic leaf $\{Q=P=0\}$, to which the coisotropic normal form of Weinstein applies \cite{weinstein,gotay82}; see also \cite{ht}, Chapter~1.
\end{proof}

The chart separates three kinds of canonical pair:
\begin{enumerate}
\item[(a)] \emph{gauge pairs} $(q^a,p_a)$: the momentum is constrained to zero and the coordinate is unconstrained --- it is the parameter along the gauge orbit, and by~\eqref{eq:HT} its velocity $\dot q^a=u^a+\partial H_c/\partial p_a$ contains an arbitrary function $u^a(t)$;
\item[(b)] \emph{second-class pairs} $(Q^\alpha,P_\alpha)$: \emph{both} members are set to zero, and the pair is invariant under the exchange $(Q,P)\mapsto(P,-Q)$;
\item[(c)] \emph{physical pairs} $(y^A,\eta_A)$, with $n-k-m$ coordinates and as many momenta.
\end{enumerate}
Type~(a) is the integrable system of Section~3 with the frequency left undetermined: the constrained momentum $p_a\weq0$ plays the part of an action with a prescribed value, and its conjugate evolves at an arbitrary rate. Type~(b) is the symmetric case.

\begin{remark}[how absolute is the asymmetry?]
In purely symplectic terms the exchange $(\tilde q,\tilde p)=(p,-q)$ rewrites $\Sigma$ as $\{\tilde q^a=0\}$, which is again coisotropic. The normal form therefore does \emph{not} by itself distinguish $q$ from $p$. The distinction comes from the polarisation inherited from the Lagrangian: by Proposition~\ref{prop:vertical} the primary constraints are vertical, and the gauge transformations of the Lagrangian fields (shifts of $A_\mu$, reparametrisations of the worldline, \dots) are shifts of coordinates. The asymmetry is a property of the \emph{pair} (symplectic manifold, polarisation), not of the symplectic manifold alone. This is exactly why (F) is true at the level of kinematics and false for the mechanical system.
\end{remark}

\begin{remark}[a mechanical Gribov argument]
Suppose that a first-class generator $\phi$ has periodic flow, as for actions. A gauge condition $\chi\weq0$ is admissible if the Faddeev--Popov factor $\PB\chi\phi=\dd\chi/\dd s$ does not vanish where $\chi=0$. Restricted to a closed orbit, $\chi$ is a smooth function on a circle; its transversal zeros come with alternating signs of $\dd\chi/\dd s$ and hence in even number. So any such condition that meets the orbit meets it at least twice, with opposite signs of the Faddeev--Popov factor. This is the same obstruction as in Theorem~\ref{thm:noconj}.
\end{remark}

\subsection{Examples}

\begin{example}[a rigid link via a Lagrange multiplier]\label{ex:link}
Let $L=\tfrac12\dot x^2+\tfrac12\dot z^2-y\,(x-z)$. Then $p_x=\dot x$, $p_z=\dot z$, $p_y=0$ and $H_c=\tfrac12p_x^2+\tfrac12p_z^2+y(x-z)$. The algorithm gives the chain
\begin{align*}
\phi_1&=p_y\weq0 &&\text{(primary)}\\
\phi_2&=x-z\weq0 &&\text{from }\dot\phi_1=-(x-z)\\
\phi_3&=p_x-p_z\weq0 &&\text{from }\dot\phi_2=p_x-p_z\\
\phi_4&=y\weq0 &&\text{from }\dot\phi_3=-2y ,
\end{align*}
and $\dot\phi_4=u$ fixes the multiplier, $u\weq0$. The non-zero brackets are $\PB{\phi_1}{\phi_4}=-1$ and $\PB{\phi_2}{\phi_3}=2$, so all four constraints are second class and the reduced phase space has dimension $6-4=2$: a free particle of mass~2 at the common position. Note the alternation momentum-type, coordinate-type, momentum-type, coordinate-type: \emph{only the primary constraint is forced to be vertical} (Proposition~\ref{prop:vertical}); later constraints may involve coordinates alone. Each second-class pair, $(\phi_1,\phi_4)$ and $(\phi_2,\phi_3)$, contains one member of each type.
\end{example}

\begin{example}[Maxwell versus Proca]\label{ex:proca}
Take signature $(-+++)$ and $\mathcal L=-\tfrac14F_{\mu\nu}F^{\mu\nu}-\tfrac12\kappa^2A_\mu A^\mu$. Then $\pi^0=0$ (primary), $\pi^i=\dot A_i-\partial_iA_0$, and up to a divergence
\[
\mathcal H_c=\tfrac12\pi^i\pi^i+\tfrac14F_{ij}F_{ij}+\tfrac12\kappa^2A_iA_i-A_0\,\partial_i\pi^i-\tfrac12\kappa^2A_0^2 .
\]
Consistency of $\pi^0\weq0$ gives $\dot\pi^0=\partial_i\pi^i+\kappa^2A_0$, i.e.\ the secondary constraint
\[
\chi=\partial_i\pi^i+\kappa^2A_0\weq0,\qquad \PB{\pi^0(\mathbf x)}{\chi(\mathbf y)}=-\kappa^2\,\delta^3(\mathbf x-\mathbf y).
\]
\emph{Maxwell} ($\kappa=0$): $\pi^0$ and $\partial_i\pi^i$ are both functions of momenta only and commute, so they are first class; $G[\epsilon]=\int\!\dd^3x\,(\dot\epsilon\,\pi^0-\epsilon\,\partial_i\pi^i)$ generates $\delta A_\mu=\partial_\mu\epsilon$, a pure shift of coordinates at fixed momenta. Gauge conditions ($A_0=0$, $\partial_iA_i=0$) are functions of coordinates. The count per point is $8-2\cdot2=4$, two polarisations.
\emph{Proca} ($\kappa\neq0$): the secondary constraint contains the \emph{coordinate} $A_0$ and pairs with $\pi^0$ into a second-class pair, a $(q,p)$ pair of type~(b). The count is $8-2=6$, three polarisations.
The mass term converts a gauge pair (asymmetric) into a second-class pair (symmetric).
\end{example}

\begin{example}[the relativistic particle and the Stueckelberg form]\label{ex:rel}
With signature $(-+++)$ and $c=1$ let $L=-m\sqrt{-\dot x^\mu\dot x_\mu}$, where the dot is the derivative with respect to an arbitrary parameter $\tau$. The Hessian has the null vector $\dot x^\nu$ (as it must for a Lagrangian that is homogeneous of degree one in the velocities), so $\rank W=3$ and there is one primary constraint, the mass shell
\[
\phi=p^\mu p_\mu+m^2\weq0,\qquad H_c=0,\qquad H_T=\tfrac12u\,\phi,\quad \dot x^\mu=u\,p^\mu,\ \dot p_\mu=0 .
\]
It is first class and a function of momenta only, as Proposition~\ref{prop:vertical} requires; it generates $x^\mu\mapsto x^\mu+\epsilon\,p^\mu$ at fixed $p$, i.e.\ reparametrisations of the worldline. A gauge condition such as $\chi=x^0-\tau$ is a function of coordinates, with Faddeev--Popov factor $\PB\chi\phi=p^0\neq0$ on the mass shell.
Compare the Stueckelberg form with an evolution parameter, $L_S=\tfrac12M\dot x^\mu\dot x_\mu$ \cite{stueck,hp73}. Now $W=M\eta_{\mu\nu}$ is non-degenerate, there are no constraints, $H=p^2/2M$, and $p^2$ is conserved without being fixed. The same worldlines are described, but in the singular formulation the mass is a \emph{constrained} value of a momentum, and in the regular formulation it is a \emph{conserved} one --- the two faces of the generator-versus-parameter asymmetry.
\end{example}

%=====================================================================
\section{The first-order action: where the equality is exact}\label{sec:first-order}
%=====================================================================

It would be wrong to conclude that (F) has no content. Consider the phase-space action
\begin{equation}\label{eq:first-order}
S[q,p]=\int\bigl(p\,\dot q-H(q,p)\bigr)\,\dd t,\qquad
S'[q,p]=-\int\bigl(q\,\dot p+H(q,p)\bigr)\,\dd t .
\end{equation}
Their difference is the boundary term $S-S'=[pq]$; both give Hamilton's equations, with
\[
\delta S=\int\bigl[\delta p\,(\dot q-H_p)-\delta q\,(\dot p+H_q)\bigr]\dd t+[p\,\delta q],\qquad
\delta S'=\int\bigl[\ \dots\ \bigr]\dd t-[q\,\delta p].
\]
Thus, as far as the \emph{equations of motion} are concerned, the choice of which variable multiplies the velocity is a total derivative. It matters only for the boundary conditions: $S$ is stationary when $q$ is fixed at the end points, $S'$ when $p$ is.

Regarded as a Lagrangian on the $2n$-dimensional configuration space $(q,p)$, $S$ is singular, because it is linear in the velocities. With momenta $\pi_q,\pi_p$ the Dirac analysis is immediate: $\phi_1=\pi_q-p\weq0$ and $\phi_2=\pi_p\weq0$, $H_c=H$, and $\PB{\phi_1}{\phi_2}=-1$, so the two constraints are second class and no further ones arise. The Dirac bracket is
\[
\PB qp_D=\PB qp-\PB q{\phi_\alpha}(C^{-1})^{\alpha\beta}\PB{\phi_\beta}p=1 ,
\]
so the Dirac bracket reproduces the canonical Poisson bracket \cite{fj88}. For $S'$ the constraints are $\pi_p+q\weq0$ and $\pi_q\weq0$: the two roles are swapped. The second-class pair $(\phi_1,\phi_2)$ is of type~(b) of Proposition~\ref{prop:normal}, so (F) is exactly the statement about second-class pairs.

For $(\theta,I)$ the first-order form is $S=\int(I\dot\theta-H(I))\,\dd t$. The action $I$ appears without a velocity and is a Lagrange multiplier enforcing $\dot\theta=\nu(I)$, whose own Euler--Lagrange equation is $\dot I=0$. If $\det H_{II}\neq0$ it can be eliminated and one obtains the second-order Lagrangian $\mathcal L(\dot\theta)$ of Section~\ref{sec:kolmogorov}; if $\det H_{II}=0$ it cannot. This closes the circle between Sections~3 and~4.

%=====================================================================
\section{Synthesis}
%=====================================================================

Table~\ref{tab:three} puts the three manifestations side by side. In each column a \emph{momentum-like} object is distinguished; it is conserved, or constrained to a fixed value, and it generates a one-parameter group that shifts a \emph{conjugate coordinate-like} object.

\begin{table}[ht]
\centering
\small
\setlength{\tabcolsep}{4pt}
\renewcommand{\arraystretch}{1.25}
\begin{tabular}{@{}>{\raggedright\arraybackslash}p{2.5cm}>{\raggedright\arraybackslash}p{4.1cm}>{\raggedright\arraybackslash}p{4.1cm}>{\raggedright\arraybackslash}p{4.1cm}@{}}
\toprule
 & \textbf{Action--angle} & \textbf{First-class constraint} & \textbf{Singular $L$ (primary)}\\
\midrule
Distinguished momentum & action $I_k$ & generator $\phi_a$ & combination $\phi_m(q,p)$ built from momenta\\
Its value & constant, $\dot I_k=0$ & weakly zero & fixed by $\FL$ on $\Sigma_1$\\
Conjugate coordinate & angle $\theta_k$ & gauge parameter $q^a$ & $q^m$ along the null directions of $W$\\
Its evolution & $\dot\theta_k=\nu_k(I)$, fixed by $H$ & $\dot q^a=u^a+\partial_{p_a}H_c$, $u^a$ arbitrary & $\dot q^m$ not determined by $L$ (arbitrary if first class)\\
Obstruction to exchange & periodic orbit: no global conjugate (Thm.~\ref{thm:noconj}) & compact orbit: no global gauge condition without copies & $\pi(\Sigma_1)=Q$: no constraint on $q$ alone (Prop.~\ref{prop:vertical})\\
Where symmetry survives & locally (universal cover) & second-class pairs & first-order action\\
\bottomrule
\end{tabular}
\caption{Three manifestations of one asymmetry.}\label{tab:three}
\end{table}

\paragraph{Corrected statement.}
The following replaces (F).
\begin{quote}
\emph{The symplectic form is symmetric under the exchange of $q$ and $p$, and so are second-class constraint pairs and the first-order action. A mechanical system carries in addition (a) the Legendre fibration $T^*Q\to Q$, with respect to which only $\mathrm{Diff}(Q)\ltimes Z^1(Q)$ is admissible; (b) a Hamiltonian whose invariant tori are compact, so that conserved ``momenta'' (actions) have non-conserved, periodic conjugates; and (c) first-class constraints, which are momentum-type generators whose conjugates are pure gauge. Relative to this structure momenta are generators --- conserved or constrained --- and coordinates are parameters along their orbits.}
\end{quote}
The invariant dichotomy is therefore ``generator versus parameter'', which is defined by a group action and not by the symplectic form alone. For a particle on $\R^n$ without constraints, (F) is harmless as a statement about the Poisson algebra and the Fourier duality. Even there, however, the Hamiltonian decides which variables are constants of the motion (Examples~\ref{ex:free} and~\ref{ex:osc}), and whether the exchanged description is still Lagrangian. The cost of using (F) as a general principle is paid whenever the phase space is the cotangent bundle of something other than $\R^n$, whenever the Lagrangian is singular, or whenever one asks which variables are conserved.

%=====================================================================
\begin{thebibliography}{99}
\bibitem{arnold} V.\,I.~Arnold, \emph{Mathematical Methods of Classical Mechanics}, 2nd ed., Springer, New York, 1989.
\bibitem{am} R.~Abraham and J.\,E.~Marsden, \emph{Foundations of Mechanics}, 2nd ed., Benjamin/Cummings, Reading, 1978.
\bibitem{dirac50} P.\,A.\,M.~Dirac, Generalized Hamiltonian dynamics, \emph{Can.\ J.\ Math.} \textbf{2} (1950) 129--148.
\bibitem{dirac64} P.\,A.\,M.~Dirac, \emph{Lectures on Quantum Mechanics}, Belfer Graduate School of Science, Yeshiva University, New York, 1964.
\bibitem{dirac27} P.\,A.\,M.~Dirac, The quantum theory of the emission and absorption of radiation, \emph{Proc.\ R.\ Soc.\ Lond.\ A} \textbf{114} (1927) 243--265.
\bibitem{sg64} L.~Susskind and J.~Glogower, Quantum mechanical phase and time operator, \emph{Physics} \textbf{1} (1964) 49--61.
\bibitem{gnh} M.\,J.~Gotay, J.\,M.~Nester and G.~Hinds, Presymplectic manifolds and the Dirac--Bergmann theory of constraints, \emph{J.\ Math.\ Phys.} \textbf{19} (1978) 2388--2399.
\bibitem{ht} M.~Henneaux and C.~Teitelboim, \emph{Quantization of Gauge Systems}, Princeton University Press, Princeton, 1992.
\bibitem{weinstein} A.~Weinstein, Symplectic manifolds and their Lagrangian submanifolds, \emph{Adv.\ Math.} \textbf{6} (1971) 329--346.
\bibitem{gotay82} M.\,J.~Gotay, On coisotropic imbeddings of presymplectic manifolds, \emph{Proc.\ Amer.\ Math.\ Soc.} \textbf{84} (1982) 111--114.
\bibitem{fj88} L.~Faddeev and R.~Jackiw, Hamiltonian reduction of unconstrained and constrained systems, \emph{Phys.\ Rev.\ Lett.} \textbf{60} (1988) 1692--1694.
\bibitem{stueck} E.\,C.\,G.~Stueckelberg, \emph{Helv.\ Phys.\ Acta} \textbf{14} (1941) 322, 588.
\bibitem{hp73} L.\,P.~Horwitz and C.~Piron, Relativistic dynamics, \emph{Helv.\ Phys.\ Acta} \textbf{46} (1973) 316--326.
\end{thebibliography}

\end{document}
